\section{Preliminaries}
\label{sec:preliminaries}

The analysis uses three related levels of variational information.  Limiting
subdifferentials formulate local stationarity and the KKT system;
level-proximal subdifferentials encode the global quadratic support required
by an exact SDG value; and Moreau regularization turns weak convexity into
well-posed proximal constructions.  This section introduces the objects in
that order.

\subsection{Variational framework}

\subsubsection{Affine geometry and limiting stationarity}

We work in finite-dimensional Euclidean spaces.  The symbols
\(\ip{\cdot}{\cdot}\) and \(\norm{\cdot}\) denote the Euclidean inner product
and norm; for a matrix, \(\norm{\cdot}\) is the induced operator norm, whereas
\(\norm{\cdot}_{F}\) is the Frobenius norm.  We write
\[
  B_r(z):=\{u:\norm{u-z}\le r\},
  \qquad
  [t]_+:=\max\{t,0\},
  \qquad
  \dist(z,D):=\inf_{u\in D}\norm{z-u}.
\]
If \(D\) is nonempty, closed, and convex, its Euclidean projection is the
single-valued map
\[
  \Pi_D(z):=\argmin_{u\in D}\norm{u-z}.
\]
Existence, uniqueness, and nonexpansiveness of this projection are standard
\citep[Example~2.25 and Corollary~12.20]{RockafellarWets1998}.

For an extended-real-valued function \(h:\R^n\to\R\cup\{+\infty\}\), its
\emph{effective domain} is
\(\operatorname{dom}h:=\{x:h(x)<+\infty\}\).  The function is proper when
this set is nonempty, and it is lower semicontinuous at \(x\) when
\(h(x)\le\liminf_{u\to x}h(u)\).  A set \(D\) is encoded by its indicator
\[
  \ind_D(x):=
  \begin{cases}
    0,       &x\in D,\\
    +\infty, &x\notin D.
  \end{cases}
\]
Thus minimizing \(h\) over \(D\) is the same as minimizing
\(h+\ind_D\) over \(\R^n\).  For a convex set \(D\),
\(\operatorname{ri}D\) denotes its interior relative to its affine hull.
These conventions follow
\citet[Sections~1.A--B and~2.H]{RockafellarWets1998}.

Recall the affine feasible set and the standing nonemptiness assumption:
\[
  \Cset:=\{x\in\R^n:Ax=b\},
  \qquad
  \Cset\cap\operatorname{dom} f\ne\varnothing .
\]
Its tangent space is \(\Tspace:=\ker A\), and its normal space is
\(\Tspace^\perp=\ran(A^\top)\).  The ordinary Lagrangian is
\begin{equation}
  \Lagr(x,y):=f(x)+\ip{y}{Ax-b},
  \qquad y\in\R^m .
  \label{eq:lagrangian}
\end{equation}

To specify the nonsmooth first-order condition precisely, let \(h\) be finite
at \(x\).  Its regular (Fr\'echet) subdifferential is specified, for
\(v\in\R^n\), by
\begin{equation}
  v\in\widehat{\partial}h(x)
  \quad\Longleftrightarrow\quad
  \liminf_{\substack{u\to x\\u\ne x}}
  \frac{h(u)-h(x)-\ip{v}{u-x}}{\norm{u-x}}
  \ge0 .
  \label{eq:frechet-subdifferential}
\end{equation}
The limiting subdifferential is characterized by
\begin{equation}
  \begin{aligned}
  v\in\partial h(x)
  \quad\Longleftrightarrow\quad&
  \text{there exist sequences }(x_k),(v_k)\text{ such that}\\
  &x_k\to x,\quad h(x_k)\to h(x),\quad
  v_k\in\widehat{\partial}h(x_k),\quad v_k\to v .
  \end{aligned}
  \label{eq:limiting-subdifferential}
\end{equation}
Both subdifferentials are empty outside \(\operatorname{dom}h\).  These
definitions
follow \citet[Definition~8.3]{RockafellarWets1998}; throughout,
\(\partial h\) without an accent means the limiting subdifferential.  When
\(h\) is continuously differentiable on a neighborhood of \(x\), the regular
and limiting subdifferentials reduce to the ordinary gradient:
\[
  \widehat{\partial}h(x)=\partial h(x)=\{\nabla h(x)\}.
\]
See \citet[Exercise~8.8(b)]{RockafellarWets1998}.
For a proper convex function, both subdifferentials coincide with the convex
subdifferential: \(v\in\partial h(x)\) if and only if
\(h(u)\ge h(x)+\ip{v}{u-x}\) for every \(u\)
\citep[Proposition~8.12]{RockafellarWets1998}.
If \(s\) is continuously differentiable on a neighborhood of \(x\), then
the limiting subdifferential obeys the exact smooth-perturbation rule
\citep[Exercise~8.8(c)]{RockafellarWets1998}:
\[
  \partial(h+s)(x)=\partial h(x)+\nabla s(x).
\]

For a closed set \(D\), the limiting normal cone is
\[
  N_D(x):=\partial\ind_D(x).
\]
For convex \(D\), this agrees with the familiar support description
\[
  N_D(x)
  =
  \{v:\ip{v}{u-x}\le0\ \text{for every }u\in D\},
  \qquad x\in D,
\]
and is empty when \(x\notin D\); see
\citet[Theorem~6.9 and Exercise~8.14]{RockafellarWets1998}.  In particular,
\(N_{\Cset}(x)=\ran(A^\top)\) at every feasible point.  The limiting Fermat
rule states that a local minimizer of a proper lower-semicontinuous function
satisfies \(0\in\partial h(x)\); for a proper convex function this condition is
also sufficient for global minimality
\citep[Theorem~10.1]{RockafellarWets1998}.

\begin{definition}[KKT pair]
  A pair \((x,y)\in\R^n\times\R^m\) is a KKT pair for
  \cref{eq:main-problem} if
  \begin{equation}
    Ax=b,
    \qquad
    0\in\partial f(x)+A^\top y .
    \label{eq:kkt}
  \end{equation}
\end{definition}

In a convex problem, \cref{eq:kkt} is sufficient for global primal--dual
optimality by the global subgradient inequality.  In a nonconvex
problem it is only a first-order condition: it can hold at a local minimum,
saddle, or maximum.

\subsubsection{Weak convexity and level-proximal support}

The central curvature hypothesis is weak convexity, in the standard sense of
\citet[Section~2]{DavisDrusvyatskiy2019}.

\begin{definition}[Weak convexity]
  A proper lower-semicontinuous function \(f:\R^n\to\R\cup\{+\infty\}\) is
  \(\rho\)-weakly convex, for \(\rho\ge0\), if
  \[
    x\longmapsto f(x)+\frac{\rho}{2}\norm{x}^2
  \]
  is convex.
\end{definition}

For comparison, a proper function \(q\) is \(\mu\)-strongly convex,
\(\mu>0\), when
\(q-\mu\norm{\cdot}^{2}/2\) is convex
\citep[Exercise~12.59]{RockafellarWets1998}.  Hence adding a quadratic with
curvature strictly larger than \(\rho\) to a \(\rho\)-weakly convex function
produces a strongly convex model and, when the model is lower
semicontinuous, a unique minimizer.

For every \(x\in\operatorname{dom}f\) and \(v\in\partial f(x)\), weak
convexity gives the global lower model
\begin{equation}
  f(u)\ge f(x)+\ip{v}{u-x}
          -\frac{\rho}{2}\norm{u-x}^2 ,
  \label{eq:weak-convex-lower}
\end{equation}
for every \(u\in\R^n\); see
\citet[Lemma~2.1]{DavisDrusvyatskiy2019} and the hypoconvexity
characterization in
\citet[Exercise~12.61]{RockafellarWets1998}.  At a differentiability point,
one has \(v=\nabla f(x)\).  If \(f\in C^2(\R^n)\), weak convexity with modulus
\(\rho\) is equivalent to the uniform Hessian bound
\(\nabla^2 f(x)\succeq-\rho I\) for every \(x\).
Weak convexity has become a standard model for nonsmooth nonconvex
optimization because a sufficiently strong quadratic regularizer restores
convexity; see \citet{DavisDrusvyatskiy2019}.

Weak convexity also permits an exact affine sum rule.  Since
\(f+\rho\norm{\cdot}^{2}/2\) and \(\ind_{\Cset}\) are convex, the qualification
\[
  \operatorname{ri}(\operatorname{dom}f)\cap\Cset\ne\varnothing
\]
gives
\begin{equation}
  \partial(f+\ind_{\Cset})(x)
  =\partial f(x)+N_{\Cset}(x)
  =\partial f(x)+\ran(A^\top),\qquad x\in\Cset\cap\dom f.
  \label{eq:affine-sum-rule}
\end{equation}
This is the convex sum rule of \citet[Theorem~23.8]{Rockafellar1970},
combined with the smooth-perturbation rule above.  Consequently,
\cref{eq:kkt} is equivalent to
\(0\in\partial(f+\ind_{\Cset})(x)\) under this qualification.
The condition is automatic when \(f\) is finite everywhere.  A KKT pair
always implies constrained stationarity; the qualification is needed for
the converse existence of a multiplier.

The lower model in \cref{eq:weak-convex-lower} uses the same global
quadratic support that governs an exact nonconvex SDG value.  The appropriate
pointwise object is the \emph{level-proximal subdifferential} introduced by
\citet[Definition~3.2]{WangWang2024}.  We parameterize it by curvature rather
than proximal stepsize: \(\beta=1/\lambda\) for \(\beta>0\).  At
\(\beta=0\), we use the global affine-support endpoint identified in
\citet[Proposition~3.6(v)]{WangWang2024}.

\begin{definition}[Level-proximal subdifferential]
  \label{def:level-prox}
  Let \(h:\R^n\to\R\cup\{+\infty\}\) be proper, let
  \(x\in\operatorname{dom}h\), and let \(\beta\ge0\).  For
  \(g\in\R^n\),
  \begin{equation}
    g\in\partial_{\mathrm p}^{\beta}h(x)
    \quad\Longleftrightarrow\quad
    h(u)\ge h(x)+\ip{g}{u-x}
      -\frac{\beta}{2}\norm{u-x}^2
    \quad\text{for every }u\in\R^n .
    \label{eq:level-prox}
  \end{equation}
\end{definition}

Unlike the classical proximal subdifferential, which requires a quadratic
support inequality only locally, \cref{eq:level-prox} is global.  Increasing
\(\beta\) weakens the support inequality, and hence
\begin{equation}
  0\le\beta_1\le\beta_2
  \quad\Longrightarrow\quad
  \partial_{\mathrm p}^{\beta_1}h(x)
  \subseteq
  \partial_{\mathrm p}^{\beta_2}h(x).
  \label{eq:level-prox-monotonicity}
\end{equation}
This monotonicity is proved in
\citet[Proposition~3.6(ii)]{WangWang2024}.  The comparison with limiting
stationarity below follows from
\citet[Proposition~3.4(ii), Proposition~3.6(iv), and
Theorem~3.9(i)]{WangWang2024}; the endpoint \(\beta=0\) is covered by
\citet[Proposition~3.6(v)]{WangWang2024}, and the boundary
\(\beta=\rho\) follows from the global weak-convexity model
\citep[Lemma~2.1]{DavisDrusvyatskiy2019}.

\begin{lemma}[Level-proximal and limiting subdifferentials]
  \label{lem:level-prox-properties}
  Let \(h\) be proper and lower semicontinuous.
  \begin{enumerate}[label=(\roman*)]
    \item For every \(\beta\ge0\),
      \(\partial_{\mathrm p}^{\beta}h(x)\subseteq\partial h(x)\).
    \item If \(h\) is \(\rho\)-weakly convex, then, for every
      \(\beta\ge\rho\),
      \begin{equation}
        \partial_{\mathrm p}^{\beta}h(x)=\partial h(x).
        \label{eq:level-prox-weak-convex}
      \end{equation}
  \end{enumerate}
\end{lemma}

Thus the level-proximal condition coincides with limiting stationarity
throughout the uniform safe-curvature regime.  Below that regime it can be
strictly stronger, because it records the least global quadratic support
needed at the particular point.  The proximal representation of this support
condition is introduced next.

\subsection{Proximal and primal--dual constructions}

\subsubsection{Moreau envelopes, proximal maps, and conjugacy}

For \(\lambda>0\), the Moreau envelope and proximal mapping of
\citet[Definition~1.22]{RockafellarWets1998} are
\begin{align}
  P_\lambda(z)=\prox_{\lambda f}(z)
    &:=\argmin_{u\in\R^n}
       \left\{f(u)+\frac{1}{2\lambda}\norm{u-z}^2\right\}, \label{eq:prox}\\
  \Moreau_\lambda(z)
    &:=\inf_{u\in\R^n}
       \left\{f(u)+\frac{1}{2\lambda}\norm{u-z}^2\right\}. \label{eq:moreau}
\end{align}
In general, \(P_\lambda(z)\) may be set valued.  If \(f\) is
\(\rho\)-weakly convex and \(\lambda\rho<1\), the proximal objective is
\((\lambda^{-1}-\rho)\)-strongly convex.  The proximal mapping is then
single valued, and the Moreau envelope is continuously differentiable with
\begin{equation}
  \nabla\Moreau_\lambda(z)
    =\frac{1}{\lambda}
      \bigl(z-\prox_{\lambda f}(z)\bigr).
  \label{eq:moreau-gradient}
\end{equation}
Thus a proximal fixed point is exactly a zero of
\(\nabla\Moreau_\lambda\), and
\(\norm{\nabla\Moreau_\lambda(z)}\) is used as a smooth stationarity
measure; see \citet[Lemma~2.2]{DavisDrusvyatskiy2019}.

The level-proximal support condition has an exact proximal representation.
For a proper lower-semicontinuous function \(h\), let
\(\prox_{\lambda h}(z)\) denote the proximal argmin in
\cref{eq:prox} with \(f\) replaced by \(h\).  For every \(\lambda>0\),
\begin{equation}
  p\in\prox_{\lambda h}(z)
  \quad\Longleftrightarrow\quad
  \frac{z-p}{\lambda}
  \in\partial_{\mathrm p}^{1/\lambda}h(p),
  \label{eq:level-prox-resolvent}
\end{equation}
which is the level-proximal resolvent relation of
\citet[Proposition~3.4(ii)]{WangWang2024}.  Unlike a
first-order optimality condition, \cref{eq:level-prox-resolvent}
characterizes a global proximal minimizer even when the proximal objective is
nonconvex.

For the affine problem, it is useful to absorb feasibility into the objective:
\begin{align}
  P_\lambda^{\Cset}(z)
    &:=\argmin_{u\in\Cset}
       \left\{f(u)+\frac{1}{2\lambda}\norm{u-z}^2\right\}, \label{eq:cprox}\\
  \Moreau_\lambda^{\Cset}(z)
    &:=\inf_{u\in\Cset}
       \left\{f(u)+\frac{1}{2\lambda}\norm{u-z}^2\right\}. \label{eq:cenvelope}
\end{align}
Since \(\ind_{\Cset}\) is convex, \(f+\ind_{\Cset}\) remains
\(\rho\)-weakly convex.  Therefore the same safe condition
\(\lambda\rho<1\) yields uniqueness on \(\Cset\).  Applying standard weakly
convex Moreau calculus to \(h=f+\ind_{\Cset}\) gives the following result
\citep[Lemma~2.5]{DavisDrusvyatskiy2022}.  It can equivalently be obtained by
completing the square and using convex Moreau calculus
\citep[Theorem~2.26 and Exercise~12.61]{RockafellarWets1998}; its fixed-point
statement is the hypoconvex resolvent identity discussed by
\citet[Theorems~3.1 and~3.7]{WangWang2024}.

\begin{lemma}[Standard constrained Moreau calculus]
  \label{lem:safe-constrained-prox}
  Suppose that \(f+\ind_{\Cset}\) is proper, lower semicontinuous, and
  \(\rho\)-weakly convex, and let \(\lambda>0\) with \(\lambda\rho<1\).  Then:
  \begin{enumerate}[label=(\roman*)]
    \item \(P_\lambda^{\Cset}\) is single valued and
      \begin{equation}
        \norm{P_\lambda^{\Cset}(z)-P_\lambda^{\Cset}(w)}
        \le
        \frac{1}{1-\lambda\rho}\norm{z-w};
        \label{eq:constrained-prox-lipschitz}
      \end{equation}
    \item \(\Moreau_\lambda^{\Cset}\) is continuously differentiable, with
      \begin{equation}
        \nabla\Moreau_\lambda^{\Cset}(z)
        =
        \lambda^{-1}
        \left(z-P_\lambda^{\Cset}(z)\right);
        \label{eq:constrained-envelope-gradient}
      \end{equation}
    \item a feasible point \(x\) satisfies
      \(P_\lambda^{\Cset}(x)=x\) if and only if
      \(0\in\partial(f+\ind_{\Cset})(x)\).  If
      \(\operatorname{ri}(\operatorname{dom}f)\cap\Cset\ne\varnothing\),
      this is equivalent to the existence of a multiplier \(y\) for which
      \((x,y)\) is a KKT pair.
  \end{enumerate}
\end{lemma}

A zero constrained-envelope gradient therefore characterizes a stationary
point of the constrained problem, not necessarily a global minimizer.

The affine proximal oracle developed in \Cref{sec:affine-prox} also uses
Fenchel conjugacy.  The Fenchel conjugate of a proper function \(h\) is
\begin{equation}
  h^*(q):=\sup_{u\in\R^n}\{\ip{q}{u}-h(u)\}.
  \label{eq:fenchel-conjugate}
\end{equation}
This is the Legendre--Fenchel transform of
\citet[Section~11.A]{RockafellarWets1998}.  If \(h\) is proper, lower
semicontinuous, and \(\mu\)-strongly convex, then \(h^*\) is continuously
differentiable and \(\nabla h^*\) is \(1/\mu\)-Lipschitz
\citep[Proposition~12.60]{RockafellarWets1998}.  Applied to
\(h(u)=\lambda f(u)+\norm{u}^{2}/2\), this identity yields the unconstrained
proximal map and the smooth dual-root potential used later.

\subsubsection{The self-centered smoothed duality gap}

Quadratic regularization also turns the classical primal--dual gap into a
finite proximal construction.  For a convex--concave Lagrangian, the
unregularized gap at \((x,y)\) is
\[
  \sup_{u,v}\{\Lagr(x,v)-\Lagr(u,y)\}.
\]
It is nonnegative because \((u,v)=(x,y)\) is admissible, but it can be
infinite.  Quadratically penalizing the comparison variables gives the
self-centered SDG of \citet[Definitions~4--5]{Fercoq2023} and
\citet[Definitions~9--10]{WalwilFercoq2025}:
\begin{equation}
  \G_{\beta_x,\beta_y}(x,y)
  :=
  \sup_{\substack{u\in\R^n\\v\in\R^m}}
  \left\{
    \Lagr(x,v)-\Lagr(u,y)
    -\frac{\beta_x}{2}\norm{u-x}^2
    -\frac{\beta_y}{2}\norm{v-y}^2
  \right\},
  \label{eq:sdg-definition}
\end{equation}
for \((x,y)\in\operatorname{dom}f\times\R^m\), with
\(\beta_x,\beta_y>0\).  This parameterization uses curvature
\(\beta_x=1/\lambda\).  Consequently the safe weak-convexity condition is
\(\beta_x>\rho\), equivalently \(\lambda\rho<1\).

In convex optimization, the SDG is a computable merit function with useful
error bounds and algorithmic interpretations \citep{WalwilFercoq2025}.  Gap
functions have also been developed for nonconvex constrained problems through
weak conjugacy and augmented Lagrangians
\citep{YalcinKasimbeyli2020}.  In the present weakly convex setting, however,
the inner comparison in \cref{eq:sdg-definition} is global.  Local
stationarity alone therefore cannot characterize its zero set.  The
decomposition in the next section exposes the corresponding global support
condition and brings the level-proximal calculus of
\cref{eq:level-prox,eq:level-prox-resolvent} into the SDG analysis.
